% Preamble for the docxplus manuscript
\usepackage{amsmath}
\usepackage{amssymb}
\usepackage{amsfonts}
\usepackage{amsthm}
\usepackage{booktabs}
\usepackage{graphicx}
\usepackage{xcolor}
\usepackage{hyperref}
% Defined here, applied via pandoc's -V linkcolor/citecolor/urlcolor. A \hypersetup
% in this file would be dead code: pandoc emits its own after the -H injection.
\definecolor{docxplusred}{HTML}{B00020}
\usepackage{caption}
\usepackage{subcaption}
\usepackage{microtype}
\usepackage{geometry}
% Tighter margins: the figures are wide, and a 1in frame forced several of them
% to scale down past the point where their 7pt labels stayed legible in print.
\geometry{top=0.75in, bottom=0.85in, left=0.85in, right=0.85in}

% Theorem environments
\theoremstyle{definition}
\newtheorem{definition}{Definition}[section]
\newtheorem{proposition}{Proposition}[section]
\newtheorem{theorem}{Theorem}[section]
\newtheorem{lemma}{Lemma}[section]
\newtheorem{corollary}{Corollary}[section]
\newtheorem{remark}{Remark}[section]
\newtheorem{example}{Example}[section]

% Float placement. LaTeX's defaults are tuned for narrow single-column journals and
% will exile a wide figure to a page of its own with 80% whitespace. These bounds
% let a figure share a page with the text that discusses it, which is the point of
% placing it there.
\renewcommand{\topfraction}{0.90}
\renewcommand{\bottomfraction}{0.75}
\renewcommand{\textfraction}{0.08}
\renewcommand{\floatpagefraction}{0.75}
\setcounter{topnumber}{2}
\setcounter{bottomnumber}{2}
\setcounter{totalnumber}{3}

% Captions: smaller than body text and clearly labelled, so a caption that carries
% real explanation still reads as apparatus rather than as prose.
\captionsetup{
  font=small,
  labelfont=bf,
  labelsep=period,
  width=0.95\textwidth,
  justification=justified,
  singlelinecheck=false,
  skip=8pt
}
